\subsection{Why identity adds little over state at the death}
A plausible reading is that the death overs compress the decision space. With few
balls and wickets left, the situation constrains intent strongly --- most batters
attack and most bowlers defend regardless of who they are --- so the state features
appear to capture most of the recoverable signal. That state explains roughly
$40\times$ more skill for runs than for wickets is consistent with the mechanism we
hypothesized: scoring is where elite finishers separate, while death-over wickets
are driven by high-variance events (mishits on yorkers and slower balls under
maximum-aggression intent) that are largely unpredictable from the observed state.
The wicket head barely improves on the base rate for any model, and the
residual-history test (H2) finds no path-dependence in wicket risk --- two
consistent indications of the same limited predictability.

\subsection{The SHAP--skill dissociation}
A result worth dwelling on is the gap between attribution and skill. TreeSHAP
credits batter identity with the majority of explanatory mass, yet identity adds
no positive held-out skill and significantly hurts runs RMSE in every one of six
seasons. The random-player placebo isolates the mechanism: shuffling player
identities collapses batter SHAP mass from $51$--$64\%$ to $7$--$8\%$, so the model
genuinely keys on real players (the attribution is faithful, not an artifact of
high-variance features); but real identity is no more useful out-of-sample than the
shuffled placebo, and in fact slightly worse. The explanation is overfitting of
non-stationary form: empirical-Bayes historical death-over rates explain
training-set variance --- which is what SHAP measures --- but a player's form does
not transfer cleanly to the next season's deliveries. Practitioners who read SHAP
mass as importance, without an out-of-sample skill check, would conclude that
death-over modeling should center on player identity. The held-out evidence says
the opposite: a state-only model is better calibrated and more stable across seeds.
This is a concrete cautionary case for pairing attribution with proper out-of-sample
scoring. We are careful about the scope of the claim: this shows that \emph{our}
identity representation --- empirical-Bayes historical death-over rates --- does not
generalize, not that player identity is irrelevant in principle (see
Section~\ref{sec:limitations}).

\subsection{What the Markov violation does and does not mean}
H2 shows a genuine but narrow departure from the Markov assumption: longer
preceding dot runs lower the next ball's expected runs, beyond state, identity,
and a deterministic pressure index. The effect is on scoring suppression, not
dismissal risk. This refines the field's backbone models --- they are
well-specified for wicket dynamics at the death but slightly mis-specified for
run-scoring intensity, where recent dot balls carry information that the
instantaneous state does not. The effect is small in absolute terms (RMSE
$+0.0065$) but robust (interval well clear of zero, correct sign).

This asymmetry --- pressure shows up in runs but not wickets --- is the finding that
most complicates a blanket ``momentum is an illusion'' reading, and it is worth
being explicit that our design identifies the effect but not its mechanism. Several
non-exclusive pathways are consistent with it. \emph{Bowler--captain adaptation:}
a run of dot balls is partly a marker of a bowler executing well (a settled yorker
or wide-line plan) and of field settings that the over-level state features do not
fully capture, so the suppression may be ongoing tactical control rather than a
psychological state of the batter. \emph{Batter recalibration:} after consecutive
dots a batter may rationally lower risk for a ball or two, or conversely be forced
into a low-percentage shot --- both depress expected runs without raising the
modelled dismissal probability. \emph{Strike rotation and matchup:} dot runs are
correlated with the striker being pinned at one end against a difficult matchup. We
cannot separate these with observational per-ball data, and we caution that the
dot-run feature is plausibly confounded with bowler quality even after the identity
and pressure-index controls; a cleaner test would condition on bowler-over identity
or use a natural experiment. What we can say is narrow and defensible: conditional
on a rich state-plus-identity model, recent dot balls carry residual information
about scoring intensity, and they do not carry residual information about wicket
risk.

\subsection{Apparent momentum and the streak-selection bias}
H3 returns to the momentum question. The naive cricket momentum estimator reports a
strong ``cold hand'' at the death, and a reader applying the standard
conditional-probability comparison --- as the cricket pressure literature
implicitly does --- would be liable to treat that as real and even ``significant''
at higher $k$. We find it is largely a finite-sequence selection artifact. Our primary quantity
is deliberately a \emph{measurement}, not a hypothesis test: the $+0.121$ gap
between the naive estimator and the Miller--Sanjurjo bias is what our sample is
well powered to estimate, and it is large and nearly equal-and-opposite to the
naive $-0.114$ ``cold hand.'' Death-over stays are short, which is the regime where
the bias is largest, so omitting the correction is particularly consequential here.
The downstream significance test is a secondary, deliberately conservative check:
the corrected estimate is near zero and no cell survives a state-stratified null
with multiple-testing control. We are careful not to overload that test --- with
$n=1{,}760$ sequences in the primary cell it is underpowered to detect a small true
momentum even if one exists, so we read it as an absence of evidence, not evidence
of absence. The defensible claim is the measured size of the mirage; the
non-significant residual is consistent with, but does not by itself prove, a true
null.

\subsection{Limitations}
\label{sec:limitations}
Several limitations bound these claims. The most important is that our finding is
about a \emph{representation} of identity, not identity in principle. We encode
players by empirical-Bayes historical death-over rates (strike rate, boundary and
dot fractions, dismissal/economy rates); a reviewer is right to note that this omits
handedness, explicit batter--bowler matchup terms, venue and conditions, and
short-horizon recent form. It is therefore possible that identity ``did not fail''
so much as our encoding of it did, and that a richer or more interaction-aware
representation would recover out-of-sample skill. The placebo shows our encoding
carries genuine, faithfully-attributed signal that nonetheless does not generalize;
it does not show that no identity representation could. Second, the identity encoders
use a simplified two-pool historical scheme (pre-$(T{-}3)$ for training,
pre-$(T{-}2)$ for evaluation) rather than strict per-match historical encoding; the
leakage audit and placebo argue against contamination, but a stricter encoding could
still shift the estimates. Third, the main calibration tables use a single held-out
test season (2026) with five seeds governing only model and resampling randomness;
the rolling-window study (Section~\ref{sec:results-robust}) extends the central
identity result across 2021--2026, but the H2 and H3 analyses are reported on 2026
only. Fourth, the published H3 permutation null used $B = 2000$ rather than the
preregistered $5000$ for runtime; the released code is set back to the
preregistered $5000$, and because the permutation count affects only $p$-value
resolution --- not the $+0.121$ bias measurement that is our primary quantity ---
the conclusions are invariant to it. Fifth, the dot-run pressure
feature is plausibly confounded with bowler quality even after our controls, so the
H2 effect should be read as associational. Finally, WASP's canonical description is
non-archival, so we cite it by its dynamic-programming formulation rather than a
peer-reviewed article.

\subsection{Future work}
Three concrete directions follow. First, build a richer identity representation ---
handedness, explicit batter--bowler matchup encodings, venue and recent-form
features --- and test whether it recovers the out-of-sample skill that the
historical-rate encoding does not; this is the most direct way to sharpen the
``state versus identity'' claim. Second, extend the rolling-window protocol to the
H2 and H3 analyses (currently reported on 2026 only) and pair it with a stricter
per-match identity encoding and a GLM-versus-LightGBM learner check, to confirm the
dissociation and the momentum null are not learner- or split-specific. Third, extend
the Miller--Sanjurjo analysis to women's T20 and to other phases (powerplay, middle
overs), where sequence lengths and base rates differ, to map how the size of the
apparent-momentum illusion varies with the regime.
